\documentclass[11pt,a4paper]{article}
\usepackage{fontspec}
\setmainfont{TeX Gyre Pagella}
\newfontfamily\headfont{Inter}[UprightFont=*-ExtraBold,BoldFont=*-ExtraBold,Extension=.otf,Path=/usr/share/fonts/opentype/inter/]
\usepackage[a4paper,margin=2.3cm]{geometry}
\usepackage{graphicx,booktabs,float,caption,array,xcolor,longtable}
\usepackage{xurl}
\usepackage[hidelinks]{hyperref}
\usepackage{titlesec}
\usepackage[round]{natbib}
\usepackage{setspace}
\setstretch{1.05}
\renewcommand{\topfraction}{0.9}
\renewcommand{\bottomfraction}{0.8}
\renewcommand{\textfraction}{0.1}
\renewcommand{\floatpagefraction}{0.75}
\setlength{\parindent}{0pt}
\setlength{\parskip}{0.55em}
\titleformat{\section}{\headfont\Large}{\thesection}{0.8em}{}
\titleformat{\subsection}{\headfont\large}{\thesubsection}{0.8em}{}
\captionsetup{font=small,labelfont=bf}
\urlstyle{same}
\hypersetup{pdftitle={The Largest City on Earth},pdfauthor={Tristan Hotham}}

\begin{document}
\begin{center}
{\Large\bfseries The Largest City on Earth\par}
\vspace{0.4em}
{\large Six thousand years of the world's biggest cities, 3700 BC to AD 2050:\\Chandler, Modelski and the UN World Urbanization Prospects 2025\par}
\vspace{1.2em}
{\large Tristan Hotham\par}
\vspace{0.3em}
{\small Independent researcher\\
ORCID: 0000-0003-4931-9349\\
DOI: \href{https://doi.org/{{meta.doi}}}{https://doi.org/{{meta.doi}}}\par}
\vspace{0.8em}
2026
\end{center}

\begin{quote}
\small
\begin{center}\textbf{Abstract}\end{center}
I trace the largest city in the world from 3700 BC to the United Nations' projections for 2050, combining Reba, Reitsma and Seto's (2016) geocoded versions of Chandler's (1987) and Modelski's (2003) compilations ({{source.places|int}} places and {{source.estimates|int}} population estimates after cleaning) with the UN World Urbanization Prospects 2025, which defines cities by the Degree of Urbanisation. At {{source.n_snaps|int}} snapshot dates, {{n_crown_sole|int}} different cities are the largest on their own. Uruk leads from 3500 to 2800 BC with up to {{reign.Uruk.peak|int}} people; the crown then passes among Sumerian, Egyptian and Mesopotamian cities, to Linzi, Alexandria and Rome, to Constantinople, Chang'an and Baghdad, and for most of 1100 to 1825 to Chinese capitals (Chinese in {{txt.china_n|int}} of {{txt.china_of|int}} snapshots). London leads from 1850 to 1914, roughly {{london.cross_in|raw}} to {{london.cross_out|raw}} by interpolation; New York, Tokyo and, in the UN's figures, Jakarta ({{un.jakarta2025|m1}} million in 2025) follow, and Dhaka is projected to pass Jakarta in {{un.dhaka_first|raw}} by a margin of {{un.dhaka_margin_2050_pct|f1}}\%. The largest city first reaches 100,000 in {{size.first_100k.y|yr}} and a million in {{size.first_1m.y|yr}} in Modelski's figures, but in the combined estimates never exceeds {{size.max_1100_1750|int}} between 1100 and 1750. The population-weighted centre of the {{source.topn|int}} largest places stays near Mesopotamia until about 800 BC, sits over Central Asia from AD 500 to 1825, swings to the North Atlantic between 1850 and 1970, and returns to Asia, reaching {{cog_sum.c2025.lat|int}}°N {{cog_sum.c2025.lon|int}}°E in 2025. The two historical compilations name different leaders at {{n_disagree|int}} of {{n_both|int}} comparable snapshots, and the 1975 seam between historical and DEGURBA figures changes rankings substantially; both are documented rather than smoothed.

\textbf{Keywords:} cities; urbanisation; historical demography; urban history; Chandler; Modelski; World Urbanization Prospects; Degree of Urbanisation; open data; mapping
\end{quote}

\section{Introduction}
Lists of the largest city in history are popular and unstable. They depend on whose estimates are used, on what counts as a city, and on how sparse records from different centuries are put side by side. This paper builds such a list transparently from open data. It asks four questions: which city was the largest at each moment from 3700 BC to the UN's projections for 2050, how large it was, how the geographic centre of the world's big cities moved, and how much the answers depend on the sources and definitions. It accompanies an interactive map (\url{{{meta.page}}}).

\section{Data}
\textbf{Historical compilations.} \citet{reba} digitised and geocoded two compilations: Tertius Chandler's \textit{Four Thousand Years of Urban Growth} \citep{chandler} and George Modelski's \textit{World Cities} \citep{modelski}. Their files on figshare \citep{figc,figa,figm} are wide tables of places (name, other names, country, coordinates, a certainty score from 1 to 3) with population figures by date. Chandler's file has {{source.chandler_rows|int}} places and {{source.chandler_values|int}} figures between {{source.chandler_first|yr}} and 1975, at {{source.chandler_cols|int}} distinct dates; Modelski's ancient file has {{source.ancient_rows|int}} places and {{source.ancient_values|int}} figures on a grid of {{source.ancient_cols|int}} dates from 3700 BC to AD 1000, listing places of 10,000 or more before 1000 BC and of 100,000 or more after; his modern file lists {{source.modern_rows|int}} cities of a million or more in AD 2000. I used the V2 files. Compared with V1 they trim leading and trailing spaces from {{source.v1_padded|int}} names and countries ({{source.v2_padded|int}} remain), correct ``Instanbul'' to Istanbul, assign Kiev to Ukraine rather than the Russian Federation, and merge V1's duplicate Guadalajara row and a misspelt ``Madgdeburg'' row into one row each, so that Chandler has {{source.v2_rows|int}} rows rather than {{source.v1_rows|int}}; no population figure differs. The files are Windows-1252 text; read as UTF-8 they fail, and read as Latin-1 they corrupt names such as Košice. {{source.lost_chars|int}} alternative names had already lost characters, stored as question marks.

\textbf{UN World Urbanization Prospects 2025.} Files 18 to 21 of the 2025 revision \citep{wup} give the 100 largest cities at each five-year date, the 100 largest of 2025 through time, capital cities, and the population of all {{source.un_cities|int}} cities that reach 50,000 between 1975 and 2050, annually, with population-weighted centroids and a plausibility rating. Cities are defined by the Degree of Urbanisation: contiguous 1 km\textsuperscript{2} grid cells with at least {{method.degurba_density|int}} inhabitants per km\textsuperscript{2} and at least {{method.degurba_min|int}} inhabitants in all. I checked that file 18's rankings and file 19's 2025 list agree exactly with file 21. Country outlines are from Natural Earth \citep{ne}.

\section{Methods}
\subsection{Corrections}
I made {{source.fixes|int}} documented corrections (Table~\ref{tab:fixes}). The main one is systematic: every BC figure for Modelski's cities in the Americas is an exact copy of the AD figure {{method.mirror|int}} years later (for example Teotihuacan at 100,000, 150,000, 100,000 and 100,000 in 900 to 600 BC and in AD 400 to 700), so I dropped the {{source.americas_copies|int}} copies. I dropped single figures out by about a factor of ten from the same place's neighbouring figures, figures attached to the wrong place, and {{source.zeros|int}} zero values, and moved two misplaced locations. Rows with the same name less than {{method.same_km|int}} km apart were merged.

{{TABLE_FIXES}}

\subsection{Places and snapshots}
I matched {{source.matched|int}} places across the two compilations by name or alternative name and distance (under {{method.match_km|int}} km), giving {{source.places|int}} places. Snapshots are Modelski's dates to AD 1000 plus Chandler's cross-sections in 1360, 650 and 430 BC and AD 361 and 622, then Chandler's benchmark years from 1100 to 1970, then the UN every five years from 1975 to 2050: {{source.n_snaps|int}} in all. At each snapshot each compilation gives a place its own figure for that date if it has one. Otherwise Chandler's figure is interpolated log-linearly between his two nearest dates for the place if they are at most {{method.gap1|int}} years apart (before AD 1000), {{method.gap2|int}} years (to 1800) or {{method.gap3|int}} years (after); Modelski's figures are interpolated only between adjacent dates of his grid. Nothing is extrapolated. The combined estimate is the mean of the compilations that give one. Places within {{source.tie_pct|f1}}\% of the top are treated as tied. From 1975 only the UN figures are used; the two definitions are not blended.

\subsection{Centre of gravity, counts and links}
The centre of gravity at each snapshot is the population-weighted mean of the unit vectors of the {{source.topn|int}} largest places (and, as a check, the 10 largest), projected back to the sphere. Counts of places of 100,000 and of a million or more are counts within the compilations before 1975 and are therefore lower bounds. I linked {{links.n|int}} historical places to the UN city containing them (same name within {{method.link_km|int}} km), which allows a place's history to run across the seam.

\section{Results}
\subsection{The succession}
Figure~\ref{fig:crown} and Table~\ref{tab:reigns} give the succession. Modelski's first date, 3700 BC, lists only Eridu ({{S.-3700.pop|int}}). Uruk leads from 3500 to 2800 BC, reaching {{reign.Uruk.peak|int}}, and is not beaten until {{txt.uruk_passed_by|raw}} in {{txt.uruk_passed_y|yr}}. Lagash, Mari and Girsu follow; Ur is the first place of 100,000 in {{size.first_100k.y|yr}}. In 2000 BC Chandler's Memphis ({{S.-2000.pop|int}}) leads because the compilations disagree on Ur ({{txt.ur2000_c|int}} in Chandler, {{txt.ur2000_m|int}} in Modelski). Egyptian capitals lead in {{txt.egypt_n|int}} of the {{txt.egypt_of|int}} snapshots between 1600 and 1000 BC; Yin leads in 1300 BC. From 1000 BC Modelski's threshold of 100,000 produces ties (Babylon, Memphis and Thebes in 900 and 800 BC; {{txt.tie700_n|int}} places in 700 BC). Linzi leads in 500 and 300 BC, Alexandria in 200 and 100 BC (Modelski's {{peak.alexandria.m_peak|int}} in 100 BC is the first million), and Rome from AD 1 to 400. Constantinople (500 and 600), Ctesiphon (622), Chang'an (700), Baghdad (tied with Chang'an in 800; alone in 900 and 1000), Kaifeng (1100), Merv (1150), Hangzhou (1200 to 1350), Nanjing (1400) and Beijing (1450 to 1600; 1750 to 1825) follow, with Istanbul in 1650 and 1700. London leads from 1850 to 1914, ties with New York in 1925, and New York leads in 1950. Tokyo leads from Chandler's 1970 figure through the UN's 1975 to 2005; Jakarta ties with Tokyo in 2010 and leads from 2015 to 2045; Dhaka leads in 2050.

\begin{figure}[tbp]\centering\includegraphics[width=\textwidth]{f_crown.png}
\caption{Population of the largest place at each snapshot (log scale), the largest in each compilation alone (dashed), and the region of the leader (band). The UN figures begin in 1975.}\label{fig:crown}\end{figure}

{{TABLE_REIGNS}}

\begin{figure}[tbp]\centering\includegraphics[width=\textwidth]{f_map.png}
\caption{Every place that has been the largest on its own, coloured by region, and the path of the population-weighted centre of the {{source.topn|int}} largest places, coloured by date.}\label{fig:map}\end{figure}

\subsection{Size}
The largest place first passes a million in {{size.first_1m.y|yr}} in Modelski's figures; in Chandler's, at the snapshot dates, not until Beijing in {{size.first_1m_c.y|raw}} ({{size.first_1m_c.pop|int}}), although his single figure for Baghdad in AD {{baghdad_peak_c.y|raw}} is {{baghdad_peak_c.pop|int}}. Between 1100 and 1750 the combined estimate for the largest place never exceeds {{size.max_1100_1750|int}}. London reaches {{size.y1900|int}} in 1900, New York {{size.first_10m.pop|int}} in 1950 and Tokyo {{size.first_20m.pop|int}} in Chandler's 1970 figure. The UN gives Tokyo {{size.y1975|int}} in 1975 and Jakarta {{size.y2025|int}} in 2025, and projects Dhaka at {{size.y2050|int}} in 2050. The compilations hold {{n.k100_1800|int}} places of 100,000 or more and {{n.m1_1800|int}} of a million in 1800, and {{n.m1_1900|int}} million cities in 1900; the UN counts {{n.m1_1975|int}} in 1975, {{n.m1_2025|int}} in 2025 and {{n.m1_2050|int}} in 2050 (Figure~\ref{fig:counts}).

\begin{figure}[tbp]\centering\includegraphics[width=\textwidth]{f_counts.png}
\caption{Places of 100,000 and of a million or more at each snapshot (log scale). Before 1975, counts within the compilations (lower bounds).}\label{fig:counts}\end{figure}

\subsection{Where the big cities were}
Figures~\ref{fig:regions} and \ref{fig:cog} show the regional make-up of the {{source.topn|int}} largest places and the longitude of their centre. Until 800 BC the centre stays between {{txt.meso_lon_min|int}}°E and {{txt.meso_lon_max|int}}°E. It moves east to {{cog_sum.c_m500.lon|int}}°E by 500 BC as Chinese and Indian cities enter the lists, back to {{cog_sum.c100.lon|int}}°E in AD 100 under Rome's weight, and stays between {{txt.lon500_1825_min|int}}°E and {{txt.lon500_1825_max|int}}°E from AD 500 to 1825. In 1900 {{reg_sum.eu_1900_n|int}} of the {{source.topn|int}} largest places are European ({{reg_sum.eu_1900|int}}\% of their population) and the centre lies at {{cog_sum.c1900.lat|int}}°N {{txt.c1900_lon|f1}}°{{txt.c1900_hem|raw}}, {{cog_sum.km_1800_1900|int}} km from its 1800 position; its westernmost point is in {{cog_sum.west_y|raw}} at {{cog_sum.west.lat|int}}°N. By 2025 it is at {{cog_sum.c2025.lat|int}}°N {{cog_sum.c2025.lon|int}}°E and by 2050 at {{cog_sum.c2050.lat|int}}°N {{cog_sum.c2050.lon|int}}°E; {{reg_sum.asia_2025_n|int}} of the {{source.topn|int}} largest cities of 2025 are in Asia and {{reg_sum.eu_2025_n|int}} in Europe. A centre on the sphere is pulled poleward when the cities are spread around a hemisphere, which explains the high latitudes of the twentieth century; part of the jump at 1975 reflects the change of source.

\begin{figure}[tbp]\centering\includegraphics[width=\textwidth]{f_regions.png}
\caption{Share of the population of the {{source.topn|int}} largest places by region at each snapshot.}\label{fig:regions}\end{figure}

\begin{figure}[tbp]\centering\includegraphics[width=\textwidth]{f_cog.png}
\caption{Longitude of the population-weighted centre of the {{source.topn|int}} (thick) and 10 (thin) largest places.}\label{fig:cog}\end{figure}

\subsection{London}
London is {{txt.london_rank1600_ord|raw}} in 1600 ({{london.p1600|int}}), {{txt.london_rank1700_ord|raw}} in 1700, {{txt.london_rank1750_ord|raw}} in 1750 and second to Beijing in 1800 ({{london.p1800|int}} against {{london.bj1800|int}}). In 1825 it is {{london.gap1825_pct|f1}}\% behind Beijing; in 1850 it leads with {{london.p1850|int}}, {{london.x1850|f2}} times Beijing. In 1925 New York ({{london.ny1925|int}}) and London ({{london.p1925|int}}) differ by {{london.gap1925_pct|f1}}\%. Linear interpolation between snapshots places London's lead at about {{london.cross_in|raw}} to {{london.cross_out|raw}} ({{london.years_est|int}} years). Under the UN definition London has {{london.un2025|int}} people in 2025 and ranks {{txt.london_rank2025_ord|raw}} (Figure~\ref{fig:london}).

\begin{figure}[tbp]\centering\includegraphics[width=\textwidth]{f_london.png}
\caption{London and its rivals, 1500 to 2050 (log scale): Chandler to 1970, the UN from 1975.}\label{fig:london}\end{figure}

\subsection{Agreement between the compilations}
At {{n_both|int}} snapshots both compilations have figures for at least {{method.min_places|int}} places; they name different leaders at {{n_disagree|int}} (Table~\ref{tab:dis}). In AD 1000 Modelski gives Baghdad {{baghdad1000.m|int}} and Chandler {{baghdad1000.c|int}}, making {{baghdad1000.c_lead|raw}} ({{baghdad1000.c_lead_pop|int}}) his largest. Chandler's highest figures for ancient Rome and Alexandria are {{txt.rome_c_share|int}}\% and {{txt.alex_c_share|int}}\% of Modelski's. {{cert.lead_c1|int}} of the {{cert.lead_n|int}} leaders before 1975 have certainty 1; the others are {{txt.cert_list|raw}}.

{{TABLE_DIS}}

\subsection{The 1975 seam}
Chandler's 1975 list and the UN's 1975 figures describe the same year with different definitions (Figure~\ref{fig:seam}). Tokyo is similar ({{seam_sum.tokyo_ch|int}} and {{seam_sum.tokyo_un|int}}); New York falls from {{seam_sum.ny_ch|int}} to {{seam_sum.ny_un|int}} ({{txt.ny_un_rank_ord|raw}} in the UN's ranking), and Jakarta rises from {{seam_sum.jakarta_ch|int}} ({{txt.jakarta_ch_rank_ord|raw}}) to {{seam_sum.jakarta_un|int}} ({{txt.jakarta_un_rank_ord|raw}}). Chandler lists {{n.ch1975_m1|int}} places of a million or more in 1975, the UN {{n.m1_1975|int}}; Modelski lists {{n.mod2000_m1|int}} in 2000, the UN {{n.un2000_m1|int}}. DEGURBA also identifies very large ``cities'' in densely settled rural areas, such as Hajipur in Bihar ({{un.hajipur.p2025|int}} in 2025).

\begin{figure}[tbp]\centering\includegraphics[width=\textwidth]{f_seam.png}
\caption{Chandler's {{txt.n_seam|int}} largest places in 1975 and the UN's figure for the city containing each.}\label{fig:seam}\end{figure}

\subsection{To 2050}
In the UN's annual series Jakarta passes Tokyo in {{un.jakarta_first|raw}}, Dhaka passes Tokyo in {{un.dhaka_over_tokyo|raw}} and Jakarta in {{un.dhaka_first|raw}}, reaching {{un.dhaka2050|int}} against Jakarta's {{un.jakarta2050|int}}. Tokyo peaks at {{un.tokyo_peak|int}} in {{un.tokyo_peak_y|raw}}. The largest gains to 2050 are in Dhaka, Karachi, Jakarta, Luanda and Cairo; Karachi enters the top ten and Seoul leaves it (Figure~\ref{fig:2050}). Sub-Saharan Africa has {{un.ssa_top100_2025|int}} of the {{method.top100|int}} largest cities in 2025 and {{un.ssa_top100_2050|int}} in 2050, with {{un.ssa_gain_share|int}}\% of the growth among the 2050 top hundred; {{un.n_shrink_top100|int}} of these cities shrink.

\begin{figure}[tbp]\centering\includegraphics[width=\textwidth]{f_2050.png}
\caption{The {{method.rise_n|int}} largest cities of 2050 in the UN's projection, 2025 (open) and 2050 (solid).}\label{fig:2050}\end{figure}

\section{Limitations}
First, the historical figures are estimates, often rounded and inferred from walls, house counts, supplies or later censuses; a factor of two is plausible for ancient cities, and the compilations disagree. Second, coverage is uneven: before AD 1000 the compilations hold {{thin.am_pre1000|int}} places in the Americas and {{thin.ssa_pre1000|int}} in sub-Saharan Africa, against hundreds in Europe and Asia later (Table~\ref{tab:cov}); absence is not evidence of absence. Third, Modelski's thresholds create ties and hide places below 100,000 after 1000 BC. Fourth, interpolation and averaging are choices; at the historical snapshots {{interp.top10_pct|f1}}\% of top-ten estimates involve interpolation. Fifth, the 1975 seam joins agglomeration-style historical figures to DEGURBA cities; I report it rather than adjust it. Sixth, 2030 to 2050 are projections. Seventh, the certainty score of \citet{reba} concerns identification and placement, not the population figure.

{{TABLE_COV}}

\section{Conclusion}
For most of six thousand years the largest city lay in a belt from Egypt and Mesopotamia to China. The industrial nineteenth century pulled the centre of the urban world to the North Atlantic for about a century; it has since returned to Asia, further south than ever before. The sources disagree often enough that any single list should be read as one defensible reading of uncertain evidence, which is how this one is offered.

\section*{Data and code availability}
The interactive map is at \url{{{meta.page}}}. Data and code are at \url{{{meta.repo}}} and archived with this paper at \url{https://doi.org/{{meta.doi}}}. Derived data and text are released under CC BY 4.0 and code under the MIT licence. Historical data: Reba, Reitsma and Seto (2016), CC BY 4.0. UN data: World Urbanization Prospects 2025, CC BY 3.0 IGO.

\section*{AI statement}
I wrote the text, with drafting and editing assistance from Claude (Anthropic), which also helped prepare the analysis code. I reviewed and checked all text and results.

\begin{thebibliography}{9}
\bibitem[Chandler(1987)]{chandler} Chandler, T. (1987) \textit{Four Thousand Years of Urban Growth: An Historical Census}. Lewiston: St David's University Press.
\bibitem[Modelski(2003)]{modelski} Modelski, G. (2003) \textit{World Cities: -3000 to 2000}. Washington: FAROS 2000.
\bibitem[Reba et al.(2016)]{reba} Reba, M., Reitsma, F. and Seto, K. C. (2016) Spatializing 6,000 years of global urbanization from 3700 BC to AD 2000. \textit{Scientific Data} 3, 160034. \url{https://doi.org/10.1038/sdata.2016.34}
\bibitem[Reba(2016a)]{figc} Reba, M. (2016) Chandler Population Data. figshare, CC BY 4.0. \url{https://doi.org/10.6084/m9.figshare.2059494.v3}
\bibitem[Reba(2016b)]{figa} Reba, M. (2016) Modelski Ancient Period Data. figshare, CC BY 4.0. \url{https://doi.org/10.6084/m9.figshare.2059497.v2}
\bibitem[Reba(2016c)]{figm} Reba, M. (2016) Modelski Modern Period Data. figshare, CC BY 4.0. \url{https://doi.org/10.6084/m9.figshare.2059500.v3}
\bibitem[United Nations(2025)]{wup} United Nations, Department of Economic and Social Affairs, Population Division (2025) World Urbanization Prospects: The 2025 Revision, files 18 to 21. CC BY 3.0 IGO.
\bibitem[Natural Earth(2026)]{ne} Natural Earth (2026) Admin 0 countries and land, 1:50m. Public domain. \url{https://www.naturalearthdata.com/}
\end{thebibliography}
\end{document}
